\section{Methodology}
\label{sec:method}

\subsection{Corpus}

The corpus holds \NumPapers{} documents, of which \NumSources{} carry a formal
claim and enter the classification. We include a work when three conditions
hold. It concerns learning or something a learning process produced, meaning a
model, its training data, its parameters, its gradients or its outputs. It
makes a formal claim, meaning a theorem, a definition, or a measurement
designed to be comparable across papers. And somebody's resources bear on that
formal claim, in the sense that deleting every resource restriction would
change what the work asserts.
% claim: corpus-is-a-working-bibliography

The third condition admits two bodies of work with no machine learning
content. Fine-grained cryptography enters because it is the only place where
the comparison between an adversary's work and an honest party's is a defined
object. The bounded-storage literature enters on the same grounds, as the
formal home of an adversary limited in space. We treat both as destinations,
and Section~\ref{sec:findings-tables} reports that nothing in the machine
learning half of the corpus has travelled to either.

\subsection{Codebook}

We classify each work on thirteen fields against a fixed vocabulary, so that
two readers assign the same value. Table~\ref{tab:fields} lists the fields.
Six record the attack scenario and threat model, and seven record what the
work assumes about resources. The party a work bounds is whoever the bound
binds, which is not always the party the work is about. A bound on a
distinguisher binds the auditor, a bound on a differentially private algorithm
binds the mechanism, and a work whose only bound falls on the party that
trains binds the learner.
% claim: classification-of-two-hundred-works

\begin{table*}[t]
\centering
\small
\caption{The thirteen fields of the classification and their values. We record a
work whose main claim bounds nobody as bounding nobody, which is an answer and
not a gap.}
\label{tab:fields}
\begin{tabular}{@{}lp{6.4cm}p{9.0cm}@{}}
\toprule
Field & Values & Records \\
\midrule
target & train data, parameters, outputs, training run, none & what the adversary tries to learn or affect \\
family & membership, attribute, reconstruct, distribution, extract, invert, forge, evade, evade audit, inject, poison, backdoor, multiple, none & the attack, in the standard taxonomy, extended where the corpus needed it \\
phase & train time, test time, post hoc, n/a & when the adversary acts; post hoc is a released model with no further training \\
access & hard label, scores, black box, white box, weights, gradients, channel, none & the strongest access the main claim grants \\
write & read only, corpus, model, responses, channel, none & what the adversary may change in the system's persistent state \\
frame & compute, statistical, capability, none & whether the bound is computational, information-theoretic, or on what the adversary may do \\
\midrule
party & attacker, learner, auditor, mechanism, none & which party the main claim bounds \\
unit & time, circuit size, circuit depth, queries, samples, memory, algorithm class, parameters, flops, perturbation, corruption, rounds, privacy budget, success probability, none & the unit of that bound \\
tightness & unstated, poly, quasi-poly, conditional, class, budget, exponent, relative, measured & how the work states the bound: not at all; the polynomial-time quantifier; a stated quasi-polynomial time; a bound under a named assumption; membership of an algorithm class; a numeric allowance; an exponent in the party's own resource; a relation to another party's work; a measurement with no definitional standing \\
baseline & stated, absent, n/a & whether the work states what public knowledge alone achieves against the target \\
second party & a unit, or none & the unit of a second bounded party, the honest party on a relative row \\
second notion & a unit, or none & a second, incomparable bound on the same party \\
assessed & explicit, implicit & whether the work treats the bound as its object or leaves it a side condition \\
\bottomrule
\end{tabular}
\end{table*}

We revised the vocabulary seven times. A classifier reporting a work it could
not place drove some revisions, and a reviewer finding two works of one kind
coded apart drove the rest. Two revisions added values for bounds that are not
computational, namely a perturbation radius, a corruption fraction and a
privacy budget. Without those values we would have recorded a work that bounds
its adversary non-computationally as bounding nobody. One revision separated a
second bounded party from one party bounded under two incomparable resource
notions. One split a bound stated as an exponent in a party's own resource
from a bound stated as a relation between two parties' work, after a review
found the two conflated. Most of the vocabulary gaps reported during the pass
fell in the single field recording the attack family, which
Section~\ref{sec:findings-tables} returns to.

\subsection{Procedure}

We read each work first through its catalogue entry, which records the
quantifier of its main claim and the hypothesis it rests on, and then through
the primary document wherever that entry was silent on a field. Four rules
decide the cases where a work admits more than one reading. We take whichever claim the abstract leads with. We take the constrained party to be whoever the bound
binds, and not whoever the paper is about. Where a work bounds two parties,
the computational bound decides the row and we record the other party's unit
beside it. And we record a work that constrains nobody as constraining nobody,
which is an answer and not a gap.

Every assignment carries its deciding sentence. For each work we hold a
verbatim fragment from the source that determines the constrained party, the
unit and the tightness, together with the definition, theorem or section it
came from. Where a work states no bound, we hold a claim of absence and the
places we checked before making it, because an unsupported claim of absence is
worth nothing. Of the \NumSources{} works, \QuoteOnFile{} carry a quotation
and
\QuoteAbsent{} an evidenced absence.
% claim: evidence-layer-per-row

\subsection{Verification}

We checked every quotation against the primary document by extracting its text
and matching a punctuation-independent fragment. We ran two passes. The first
asked whether a quotation appeared anywhere in the corpus, and found none
absent. We had asked a weaker question than the failure it was meant to catch. The
second pass binds each work to its own document and asks whether the quotation
appears in \emph{that} document. It found four works carrying a neighbour's
quotation, an off-by-one rotation within one batch, which the first pass had
passed. We re-extracted all four from their own sources. The second pass binds
\QuoteBound{} of the \QuoteOnFile{} quotations to their document and verifies
every one there. We could not bind the remaining \QuoteUnbound{} to a document
by title, so we checked them against the corpus only.

Eleven entries failed an automated check and all eleven were false alarms, six
from fragments spanning a line break, two from an error in the checker, one
from a claim of absence written as prose, and two from mathematical notation.
A single automated pass is not sufficient evidence that a quotation is absent,
and a pass that asks a weaker question than the failure mode is worse than
none, because it returns a clean report.

\subsection{Threats to validity}

Three limits bear on the counts of Section~\ref{sec:findings-tables}. One
reader classified each work, so the agreement rate between readers is
unmeasured, and the fixed vocabulary and decision rules stand in for that
measurement.
% claim: single-coder-agreement-unmeasured

We re-checked one field in full before reporting it. The column recording
whether a work states what an adversary achieves from public knowledge alone
was applied under two readings during the first pass, one asking for the
success rate available without access to the target and one counting any
explicit comparison quantity. We re-read every row under the narrower
definition and
\BaselineChanged{} changed value. Where two readers checked overlapping rows
they agreed on \AgreeOverlap{} of \AgreeTotal{}, which is the only
inter-reader agreement figure
available here.

The field recording the attack family absorbed most of the vocabulary gaps
reported during the pass. We built that field around attacks on a deployed
model, and the corpus also contains attacks on a cryptographic primitive, on a
communication channel, on a training distribution and on an agent's control
flow. Its categories are therefore the least reliable of the thirteen, and
Figure~\ref{fig:sankey} collapses them for that reason.
